\chapter{Introduction}

The EEG signal has been used in both research and clinical applications to classify brain states and diagnose various conditions, such as epilepsy or sleep disorders\cite{jain}.
What initially required manual inspection is being superseded by automated computer methods.
The ability to infer signals from EEG recordings has sparked hope for the eventual use of EEG in Brain-Computer Interfaces (BCI).
Our work explores one such avenue.

We analyze datasets of EEG recordings while participants listened to music, so that the dataset is a set of EEG + music pairs.
We attempt to infer music properties from the EEG signals with machine learning methods.
Expecting the problem of full music reconstruction to be hard, we ramp up to it by solving few prediction tasks attempting to decode a discrete set of features, demonstrating first that any such decoding is possible.

We start with a task of predicting a song name from the EEG signal.
We succeed at this song classification task but observe poor cross-subject generalization.
We try to explain this result with an experiment that attempts to perform song identification using a classical EEG sequence matching algorithm.

We attempt emotion classification, where emotion labels are part of the dataset, which is similar task to song identification. 
Looking for more musical features,
we also define a novel classification task based on the number of note onsets in a fragment of a musical recording and solve it using our methods.

Lastly, we tackle the harder than classification problem of music decoding from EEG in a form of a spectrogram and achieve results good enough to perform classification on, but the poor generalization problem persists.
We also try to decode the rhythm of the music with a different method, trying to predict it in a form of an envelope.

This work researches the limits, applicability and modifications of methods for EEG classification tasks and elucidates some aspects of why music decoding from EEG remains a challenging problem.

The work is structured as follows:
\begin{enumerate}
    \item in Chapter~\ref{chap:eeg-signal} we describe the EEG signal broadly,
    \item in Chapter~\ref{chap:methods} and Chapter~\ref{chap:datasets} we describe the methods and datasets used in this work,
    \item in Chapter~\ref{chap:experiments} we provide the results
    \item and in Chapter~\ref{chap:ablation} we provide ablation studies.
\end{enumerate}